\documentclass[10pt,a4paper]{amsart}
\usepackage[margin=19mm]{geometry}
\usepackage{amsmath,amssymb,mathtools}
\usepackage[scr=rsfs,cal=euler]{mathalfa}
\usepackage{xfrac}
\usepackage[unicode,hidelinks,bookmarks=false]{hyperref}
\newcommand{\ie}{i.e.\ }
\newcommand{\cf}{cf.\ }
\newcommand{\resp}{respectively\ }
\newcommand{\Subsection}{\subsection}
\providecommand{\nobreakdash}{\nobreak\hbox{-}}
\setlength{\emergencystretch}{2em}
\multlinegap0pt
\usepackage{bm}
\usepackage{bbm}
\usepackage{dsfont}
\usepackage{xspace}
\usepackage{cancel}
\usepackage{mathtools}
\usepackage{amsthm}
\makeatletter
\@ifundefined{thm}{\newtheorem{thm}{Thm}}{}
\@ifundefined{lemma}{\newtheorem{lemma}[thm]{Lemma}}{}
\makeatother
\title[Complex dynamics in two-dimensional coupling of quadratic ma]{Complex dynamics in two-dimensional coupling of quadratic maps}
\author{Anca Radulescu, Eva Kaslik, Alexandru Fikl}
\date{Source: arXiv:2303.09329v4 (CC BY 4.0)}
\begin{document}
\maketitle

\noindent\emph{Source and licence.} Complex dynamics in two-dimensional coupling of quadratic maps. Version arXiv:2303.09329v4 (CC BY 4.0), distributed under the Creative Commons Attribution 4.0 International licence, as recorded in the arXiv licence metadata for this exact version and shown on the abstract page https://arxiv.org/abs/2303.09329v4. Full rights notice: licenses/arxiv-2303.09329-CC-BY-4.0.txt.\par\smallskip

\noindent\textit{Editorial scope.} Counted unit: the Lemma whose source label is \texttt{case1} in the article (source0.tex lines 171--179, source label \texttt{case1}), delivered together with the article's own complete original proof of it (source0.tex lines 181--236, one \texttt{proof} environment of 56 lines). No mathematics is written, repaired or re-derived: statement and proof are the article's own text line for line. Required context retained uncounted: none. Every definition, display and label of those lines is retained unchanged. Omitted from this package and not redistributed: 1--170, 180--180, 237--1003. Nothing inside a retained excerpt is omitted or altered: every retained body is delivered exactly as the article's lines are written, so no omission receipt of any kind accompanies this package. Citations: the retained excerpts carry no citation command at all, so no bibliography entry of the article is transcribed and no reference is redistributed. Reference adaptation: none was needed and none was made; every reference inside the delivered unit resolves inside it, so no unresolved reference or citation is produced. Printed slips kept literal: no printed word, symbol, hypothesis or number of the counted statement or of its proof was altered. Layout and class: the article's own class is \texttt{article} with packages including iftex, xparse, fontenc, fontspec, geometry, amsmath, amssymb, amsthm, xfrac, xcolor, graphicx, hyperref, ulem, wrapfig; the wrapper is the lane's pinned \texttt{amsart} editorial wrapper at 10pt on A4, which restates the theorem-like environments the retained text needs and adds the article's own macro definitions that the retained text needs (none), with extra wrapper packages (bm, bbm, dsfont, xspace, cancel, mathtools). The delivered numbering is the unit's own. Assets: no retained span contains a figure environment, a graphics-inclusion command, a PDF-page inclusion, a table or a TikZ picture; no image file is copied into this package, the package declares no asset dependency and no image file is redistributed with it. Licence: version arXiv:2303.09329v4 of this article is distributed under the Creative Commons Attribution 4.0 International licence, as recorded in the arXiv licence metadata for this exact version and shown on the abstract page https://arxiv.org/abs/2303.09329v4. A full rights notice is delivered at licenses/arxiv-2303.09329-CC-BY-4.0.txt. Characters: every retained excerpt is pure ASCII, so the delivered engine, XeLaTeX with its default Latin Modern text font, reports no missing character.
\begin{lemma} \label{case1}
Let the connectivity matrix $A$ be such that $\Delta = |a f| - |bd| \neq 0$ and $(z_1(0), z_2(0))$ arbitrary. Then, there exist a large enough $M > 0$ and an $n \ge 0$, such that
\[
\max\{|z_1(n)|, |z_2(n)|\} > M \implies \max\{|z_1(n + 1)|, |z_2(n + 1)|\} > 2 M.
\]

In particular, this is true for the critical orbit $z_1(0) = z_2(0) = 0$ that describes the
equi-M set.
\end{lemma}

\begin{proof}
Assuming that $\Delta = |af | - |bd | \neq 0$, we define
\begin{equation} \label{case1.m1m2}
M_1 = K_1 + \sqrt{K_1 (K_1 + |c|)}
\text{ and }
M_2 = K_2 + \sqrt{K_2(K_2 + |c|)},
\end{equation}
where
\[
K_1 = \left(\frac{|b| + |f|}{|\Delta|}\right)^2
\text{ and }
K_2 = \left(\frac{|a| + |d|}{|\Delta|}\right)^2.
\]

Let $M > \max\{M_1, M_2\}$. To simplify the notation, we abbreviate $z_1(n)$ and $z_2(n)$ to $z_1$ and $z_2$ for the rest of the proof, when there is no danger of confusion. Suppose that, for some $n \geq 0$, we have $|z_1| > M$ or $|z_2| > M$. We will show that $|z_1(n + 1)| > 2 M$ or $|z_2(n + 1)| > 2 M$ using a standard proof by contradiction. Suppose that both $z_1(n + 1) \leq 2M$ and $z_2(n + 1) \leq 2M$. Then, starting from $|z_1(n + 1)| \le 2 M$, we have that
\[
2 M
\geq |z_1(n + 1)| = |(a z_1 + b z_2)^2 + c|
\geq |a z_1 + b z_2|^2 - |c|
\geq (|a z_1| - |b z_2|)^2 - |c|,
\]
by repeated use of the reverse triangle inequality. Applying the analogous steps to $|z_2(n + 1)| \le 2 M$, we obtain
\[
2 M
\geq |z_2(n + 1)| = |(d z_1 + f z_2)^2 + c|
\geq |d z_1 + f z_2|^2 - |c|
\geq (|d z_1| - |f z_2|)^2 - |c|,
\]
Calling $\xi_1=|a z_1| - |b z_2|$ and $\xi_2 = |d z_1| - |f z_2|$, for convenience, we obtain that
\begin{equation} \label{case1.proof1}
    |\xi_1| \leq \sqrt{2M + |c|}
    \quad \text{and} \quad
    |\xi_2| \leq \sqrt{2M + |c|}.
\end{equation}

Using the inequalities from \eqref{case1.proof1}, we have that
\begin{equation*}
\big||f| \xi_1 - |b| \xi_2\big|
    = \big|(|a f| - |b d|) z_1\big|
    \leq |f| |\xi_1| + |b| |\xi_2|
    \leq (|b| + |f|) \sqrt{2M + |c|}.
\end{equation*}

This result provides a bound on $z_1$, under the assumption that $\Delta \ne 0$, i.e.
\begin{equation} \label{ineq_z1}
|z_1| \leq \frac{|b| + |f|}{|\Delta|} \sqrt{2 M + |c|} = \sqrt{K_1 (2 M + |c|)}.
\end{equation}

We are looking for a condition on $M$ that implies
$\sqrt{K_1 (2 M + |c|)} < M$. Squaring both sides, this is equivalent to requiring that $h_1(M) = M^2 - K_1 (2 M +|c|) > 0$. We can guarantee $h_1$ remains positive by taking $M > M_1$, the larger of its two roots. Hence, if $M>M_1$, we have that $|z_1| \leq M$. Analogously, we obtain
\begin{equation*}
   |z_2| \leq \sqrt{K_2 (2 M + |c|)},
\end{equation*}

which in turn is smaller than $M$, $M > M_2$, where $M_2$ is the larger of the roots of $h_2(M) = M^2 - K_2 (2 M + |c|)$ (as defined in~\eqref{case1.m1m2}). Since $z_1$ and $z_2$ cannot be simultaneously smaller than $M$, the contradiction follows.
\end{proof}

\end{document}
